\section{Future Roadmap}
\label{sec:roadmap}

The evolution of NavHAL is driven by the needs of modern embedded engineering. Our roadmap focuses on several key areas of growth:

\subsection{NavHAL++ (C++ Wrapper Layer)}
To cater to the growing popularity of C++ in embedded systems, we are developing a header-only C++ wrapper layer. This will provide:
\begin{itemize}
    \item Object-oriented peripheral management.
    \item Template-based type safety for register configurations.
    \item Support for RAII (Resource Acquisition Is Initialization) for power management.
\end{itemize}

\subsection{RTOS Ecosystem Integration}
Future versions of NavHAL will include "RTOS Ports" for popular kernels like FreeRTOS, Zephyr, and RT-Thread. This includes thread-safe peripheral access and standardized OS-abstraction layers.

\subsection{Automated Hardware Testing}
We are expanding our testing suite to include "Hardware-in-the-loop" (HIL) testing, ensuring that every commit is verified on actual silicon before being merged.

\section{Conclusion}
NavHAL is more than just a library; it is a fundamental shift in how embedded software is architected. By providing a robust, high-performance, and truly portable abstraction layer, NavHAL empowers engineering teams to focus on innovation rather than hardware minutiae.

As \textbf{NAVRobotec Pvt Ltd} continues to expand the NavHAL ecosystem, it remains committed to the mission of creating seamless integration between human ingenuity and machine capability.

\vfill
\begin{center}
    \textbf{NAVRobotec Pvt Ltd} \\
    Building the Future of Embedded Systems.
\end{center}
